% GENERATED FILE. Edit canonical lessons, then run npm run generate:books.
% canonical-source-hash: fnv1a64:5bdff4f54d772657
% canonical-lessons: LA-C359-pondero, LA-C359-delibero, LA-C359-dubitatio, LA-C359-probabilis, LA-C359-verisimilis

\chapter{Weighing a Question}
\label{ch:la-weighing-a-question}
\hlchaptermodality{359}

\begin{chapteropening}
\textbf{By the end of this chapter:} \emph{I can say that I weigh a question and think it over, that I am in doubt, and that a thing is likely.}

Complete the last lesson of chapter 359: I can say that I weigh a question and think it over, that I am in doubt, and that a thing is likely.
\end{chapteropening}

\section[ponderō, ponderāre]{ponderō, ponderāre — to weigh (up), to ponder}

\label{lesson:LA-C359-pondero}



\begin{quote}
Before the new one: say the Latin for agreement, then the Latin for a dispute.
\end{quote}

\subsection*{You'll want to know: ponderō, ponderāre}
\textbf{ponderō, ponderāre} — \textquotedblleft{}to weigh (up), to ponder\textquotedblright{}.

\begin{quote}
Verba eius dīligenter ponderāvī.
\end{quote}

\begin{quote}
\emph{I weighed his words carefully.}
\end{quote}

\emph{ponderō, ponderāre, ponderāvī, ponderātum}. Built on \emph{pondus}, a weight: you weigh words as a merchant weighs goods. English \emph{ponder}.

\begin{cousinweb}
\emph{ponderō} is built on \emph{pondus, ponderis} \textquotedblleft{}a weight\textquotedblright{} (Lewis \& Short; EtymDB lists \emph{ponderō} \textquotedblleft{}to weigh\textquotedblright{}). English \emph{ponder} comes from it through French.
\end{cousinweb}

\begin{grammarlens}[title={What you've built}]
One more word for telling, explaining and describing at B1.
\end{grammarlens}

\subsection*{Your turn}
\begin{itemize}
\raggedright
  \item \emph{Say it:} \emph{ponderō, ponderāre}
  \item \emph{Say it:} \emph{ponderō, ponderāre}, once more
  \item \emph{Say it:} a sentence of your own with \emph{ponderō, ponderāre}
  \item \emph{Recall:} say the Latin for agreement, then the Latin for a dispute, then say \emph{ponderō, ponderāre} again
\end{itemize}

\begin{tcolorbox}[breakable,colback=teal!4,colframe=teal!35!black,title={Before you move on}]
What does ponderō, ponderāre mean? (\textquotedblleft{}To weigh (up), to ponder\textquotedblright{}.) Say it once more.
\end{tcolorbox}
\section[dēlīberō, dēlīberāre]{dēlīberō, dēlīberāre — to deliberate, to think over}

\label{lesson:LA-C359-delibero}



\begin{quote}
Before the new one: say the Latin for a dispute, then the Latin for to weigh.
\end{quote}

\subsection*{You'll want to know: dēlīberō, dēlīberāre}
\textbf{dēlīberō, dēlīberāre} — \textquotedblleft{}to deliberate, to think over\textquotedblright{}.

\begin{quote}
Diū dēlīberāvimus; tandem cōnstituimus.
\end{quote}

\begin{quote}
\emph{We deliberated for a long time; at last we decided.}
\end{quote}

\emph{dēlīberō, dēlīberāre, dēlīberāvī, dēlīberātum}. \emph{Dē rē dēlīberāre}, to think a matter over. English \emph{deliberate}.

\begin{cousinweb}
\emph{dēlīberō} is built from \emph{dē} and a verb made on \emph{lībra} \textquotedblleft{}a balance, a pair of scales\textquotedblright{} (EtymDB; Lewis \& Short): to weigh a decision on the scales. English \emph{deliberate} comes from it.
\end{cousinweb}

\begin{grammarlens}[title={What you've built}]
One more word for telling, explaining and describing at B1.
\end{grammarlens}

\subsection*{Your turn}
\begin{itemize}
\raggedright
  \item \emph{Say it:} \emph{dēlīberō, dēlīberāre}
  \item \emph{Say it:} \emph{dēlīberō, dēlīberāre}, once more
  \item \emph{Say it:} a sentence of your own with \emph{dēlīberō, dēlīberāre}
  \item \emph{Recall:} say the Latin for a dispute, then the Latin for to weigh, then say \emph{dēlīberō, dēlīberāre} again
\end{itemize}

\begin{tcolorbox}[breakable,colback=teal!4,colframe=teal!35!black,title={Before you move on}]
What does dēlīberō, dēlīberāre mean? (\textquotedblleft{}To deliberate, to think over\textquotedblright{}.) Say it once more.
\end{tcolorbox}
\section[dubitātiō, dubitātiōnis]{dubitātiō, dubitātiōnis — doubt, hesitation}

\label{lesson:LA-C359-dubitatio}



\begin{quote}
Before the new one: say the Latin for to weigh, then the Latin for to deliberate.
\end{quote}

\subsection*{You'll want to know: dubitātiō, dubitātiōnis}
\textbf{dubitātiō, dubitātiōnis} — \textquotedblleft{}doubt, hesitation\textquotedblright{}.

\begin{quote}
Sine ūllā dubitātiōne tibi crēdō.
\end{quote}

\begin{quote}
\emph{I believe you without any doubt.}
\end{quote}

\emph{dubitātiō, dubitātiōnis} f., the noun of \emph{dubitō}, to doubt. \emph{Sine dubitātiōne}, without hesitation. English \emph{doubt} comes through French from \emph{dubitāre}.

\begin{cousinweb}
\emph{dubitātiō} is the noun built on \emph{dubitō} \textquotedblleft{}to doubt\textquotedblright{} (EtymDB), which belongs with \emph{dubius} \textquotedblleft{}wavering between two\textquotedblright{} (Lewis \& Short).
\end{cousinweb}

\begin{grammarlens}[title={What you've built}]
One more word for telling, explaining and describing at B1.
\end{grammarlens}

\subsection*{Your turn}
\begin{itemize}
\raggedright
  \item \emph{Say it:} \emph{dubitātiō, dubitātiōnis}
  \item \emph{Say it:} \emph{dubitātiō, dubitātiōnis}, once more
  \item \emph{Say it:} a sentence of your own with \emph{dubitātiō, dubitātiōnis}
  \item \emph{Recall:} say the Latin for to weigh, then the Latin for to deliberate, then say \emph{dubitātiō, dubitātiōnis} again
\end{itemize}

\begin{tcolorbox}[breakable,colback=teal!4,colframe=teal!35!black,title={Before you move on}]
What does dubitātiō, dubitātiōnis mean? (\textquotedblleft{}Doubt, hesitation\textquotedblright{}.) Say it once more.
\end{tcolorbox}
\section[probābilis, probābile]{probābilis, probābile — likely, credible; worth approving}

\label{lesson:LA-C359-probabilis}



\begin{quote}
Before the new one: say the Latin for to deliberate, then the Latin for doubt.
\end{quote}

\subsection*{You'll want to know: probābilis, probābile}
\textbf{probābilis, probābile} — \textquotedblleft{}likely, credible; worth approving\textquotedblright{}.

\begin{quote}
Tua causa mihi probābilis vidētur.
\end{quote}

\begin{quote}
\emph{Your case seems credible to me.}
\end{quote}

\emph{probābilis, probābile}: two endings. Built on \emph{probō}, to approve, to test: a thing you can approve. For a view that is likely but not certain. English \emph{probable}.

\begin{cousinweb}
\emph{probābilis} is built on \emph{probō} \textquotedblleft{}to approve, to test\textquotedblright{} (EtymDB; Lewis \& Short), which comes from \emph{probus} \textquotedblleft{}good, upright\textquotedblright{}. English \emph{probable} comes from it.
\end{cousinweb}

\begin{grammarlens}[title={What you've built}]
One more word for telling, explaining and describing at B1.
\end{grammarlens}

\subsection*{Your turn}
\begin{itemize}
\raggedright
  \item \emph{Say it:} \emph{probābilis, probābile}
  \item \emph{Say it:} \emph{probābilis, probābile}, once more
  \item \emph{Say it:} a sentence of your own with \emph{probābilis, probābile}
  \item \emph{Recall:} say the Latin for to deliberate, then the Latin for doubt, then say \emph{probābilis, probābile} again
\end{itemize}

\begin{tcolorbox}[breakable,colback=teal!4,colframe=teal!35!black,title={Before you move on}]
What does probābilis, probābile mean? (\textquotedblleft{}Likely, credible; worth approving\textquotedblright{}.) Say it once more.
\end{tcolorbox}
\section[vērisimilis, vērisimile]{vērisimilis, vērisimile — probable, like the truth}

\label{lesson:LA-C359-verisimilis}



\begin{quote}
Before the new one: say the Latin for doubt, then the Latin for likely.
\end{quote}

\subsection*{You'll want to know: vērisimilis, vērisimile}
\textbf{vērisimilis, vērisimile} — \textquotedblleft{}probable, like the truth\textquotedblright{}.

\begin{quote}
Nōn est vērisimile eum hoc fēcisse.
\end{quote}

\begin{quote}
\emph{It is not likely that he did this.}
\end{quote}

\emph{vērisimilis, vērisimile}, also written as two words, \emph{vērī similis}, like the truth. \emph{Vērisimile est} + accusative and infinitive. English \emph{verisimilitude}.

\begin{cousinweb}
\emph{vērisimilis} grew out of the phrase \emph{vērī similis} \textquotedblleft{}like the truth\textquotedblright{} (Lewis \& Short), from \emph{vērus} \textquotedblleft{}true\textquotedblright{} and \emph{similis} \textquotedblleft{}like\textquotedblright{}. English \emph{verisimilitude} comes from it.
\end{cousinweb}

\begin{grammarlens}[title={What you've built}]
One more word for telling, explaining and describing at B1.
\end{grammarlens}

\subsection*{Your turn}
\begin{itemize}
\raggedright
  \item \emph{Say it:} \emph{vērisimilis, vērisimile}
  \item \emph{Say it:} \emph{vērisimilis, vērisimile}, once more
  \item \emph{Say it:} a sentence of your own with \emph{vērisimilis, vērisimile}
  \item \emph{Recall:} say the Latin for doubt, then the Latin for likely, then say \emph{vērisimilis, vērisimile} again
\end{itemize}

\begin{tcolorbox}[breakable,colback=teal!4,colframe=teal!35!black,title={Before you move on}]
What does vērisimilis, vērisimile mean? (\textquotedblleft{}Probable, like the truth\textquotedblright{}.) Say it once more.
\end{tcolorbox}
